\section{Scope, Chronology, and Evidence}
\label{sec:bellarmine-apparatus}

\subsection{Identity and bounds}

\textbf{Name and life.} This study uses Saint Robert Bellarmine, Italian
Roberto Francesco Romolo Bellarmino, of the Bellarmini family: Jesuit religious, priest, preacher,
teacher, controversial theologian, catechist, spiritual writer, religious
superior, Roman consultor, cardinal, archbishop of Capua, canonized saint, and
Doctor of the Universal Church.  He was born at Montepulciano on 4 October
1542 and died at Sant'Andrea al Quirinale in Rome on 17 September 1621.
``Hammer of heretics'' is an early posthumous epithet, not a lifetime office.

\textbf{Places and institutions.} The principal itinerary is Montepulciano;
Rome; Florence; Mondovì; Padua; Louvain; Paris during the French legation;
Naples; Capua; and Rome again.  Venice and England are principally fields of
written and curial controversy, not places Bellarmine is claimed to have
visited for those disputes.  Governing institutions include the Society of
Jesus and its colleges, the post-Tridentine Latin Church, the Roman Curia, the
Congregations of the Index and Holy Office, the College of Cardinals, the
archdiocese of Capua, and the papacy.

\textbf{Language, ritual, and jurisdiction.} Bellarmine wrote principally in
Latin and Italian.  Hebrew and Greek learning is stated only where directly
evidenced and is not equated with modern critical philology.  The study concerns
Latin Catholic institutions, their confessional conflicts, and later Roman
Catholic reception.  It does not infer one uniform legal force across every
early-modern state, translate indirect papal power into a modern constitution,
or treat a current Roman memorial as every local or historical calendar.

\textbf{Currentness.} Current official, shrine, institutional, and catalog
claims retained here were checked through 20 July 2026.  The current Roman
General Calendar lists an optional memorial on 17 September.  Present shrine
custody and calendar claims are mutable.  Paul~VI's 1967 act establishes
Bellarmine as equal principal patron of the Archdiocese of Capua with Saint
Stephen Protomartyr; no universal patronage of a professional or devotional
class is asserted because no competent act with defined scope was located.

\subsection{Reader, question, and thesis boundary}

The intended reader is a serious non-specialist, cleric, student, or researcher
who needs a chronological life with explicit source limits.  The governing
question is how Bellarmine's Jesuit teaching, controversy, catechesis,
censorial and inquisitorial service, cardinalate, residential pastorate, and
spiritual writing generated both durable Catholic authority and a contested
legacy.

The historical thesis is that Bellarmine made post-Tridentine doctrine
systematic, historically argued, and teachable, and that he did so through
institutions capable of pastoral service and coercive force.  His prayer,
charity, catechesis, and Capuan government belong to the record alongside his
defense of lethal punishment for heretics, concrete role in Bruno's process,
strong papal claims, and 1616 action toward Galileo.  Canonization and Doctor
status are received according to the object of the official acts; this study
does not issue those judgments or extend them to every opinion and decision.

Included are the controlled life and office sequence; principal writing
genres; representative theological, scriptural, political, catechetical, and
spiritual concerns; French, Vulgate, Index, \emph{de auxiliis}, Capuan,
Venetian, English, Bruno, and Galileo episodes; death; earliest lives; cult;
relic and image claims; beatification; canonization; Doctor status; and current
commemoration.  Excluded are a critical edition, complete doctrinal commentary,
full reconstruction of each tribunal and controversy, canonization dossier,
independent medical or relic judgment, universal patronage without an act, and
an ecclesiastical verdict on this publication.

\subsection{Evidence key}

\begin{evidencelegend}
\evidencelegendrow{A}{Subject-authored evidence: Bellarmine's autobiography,
letters, theological and controversial works, catechisms, memorials, and
spiritual books, controlled by date, genre, audience, revision, and
transmission.  The 1613 autobiography is retrospective self-interpretation.}
\evidencelegendrow{N}{Contemporary or near-contemporary external evidence:
Jesuit, curial, episcopal, diplomatic, trial, correspondence, and other records
outside Bellarmine's unilateral control.  Institutional records have legal and
adversarial purposes.}
\evidencelegendrow{E}{Early received report: the 1622 and 1624 lives, process
testimony, and early cult memory.  These may preserve close evidence while
serving edifying and cause-promoting ends.}
\evidencelegendrow{L}{Later hagiographic, devotional, artistic, quotation,
relic, patronage, and polemical tradition; evidence first for the age and
community that transmit it.}
\evidencelegendrow{M}{Official ecclesial reception, limited to the act's
object: beatification, canonization, Doctor declaration, liturgical
commemoration, papal catechesis, or a competent current custodial claim.}
\evidencelegendrow{K}{Named modern critical reconstruction, especially for
chronology, institutional motive, lost records, bibliographic history, Bruno,
Galileo, and political controversy, with premises and alternatives retained.}
\evidencelegendrow{S}{Source-grounded project synthesis; an editorial
connection supported by identified evidence and never assigned to an earlier
witness or ecclesial act.}
\end{evidencelegend}

\subsection{Detailed chronology}

Dates are Common Era.  An exact day belongs only where a controlled witness or
official act supplies it.  Approximate ranges preserve sequence without
inventing a transcript, degree, appointment act, or institutional motive.

\begin{biotimeline}
\biotimelinerow{4 Oct.~1542}{Montepulciano}{Birth to Vincenzo Bellarmino and
Cinzia Cervini.}{A/K; later official concurrence}
\biotimelinerow{1540s--1550s}{Montepulciano}{Humanist schooling, literary and
musical formation, and remembered religious development.}{A retrospective}
\biotimelinerow{1559--1560}{Montepulciano / Il Vivo}{Discernment and Jesuit
probation amid family negotiation.}{A sequence; fine dates open}
\biotimelinerow{20 Sept.~1560}{Rome}{Admission to the Society of Jesus on the
vigil of Saint Matthew.}{A/K; later official concurrence}
\biotimelinerow{1560--1563}{Rome}{Early Jesuit study at the Roman College.}{A;
exact courses unresolved}
\biotimelinerow{1563--1567}{Central / northern Italy}{Formation, teaching, and
preaching assignments including Florence and Mondovì.}{A/K approximate}
\biotimelinerow{1567--spring 1569}{Padua}{Further study before transfer to the
Low Countries.}{A/K; degree sequence open}
\biotimelinerow{spring 1569}{Louvain}{Sent to continue formation, preach, and
teach.}{A/K strong}
\biotimelinerow{25 Mar.~1570}{Louvain setting}{Priestly ordination after the
liturgical sequence remembered in the autobiography.}{A sequence / K date;
later official concurrence; exact place open}
\biotimelinerow{1570--1576}{Louvain}{Preaches and teaches theology while
continuing study and Hebrew work.}{A/K strong range}
\biotimelinerow{July 1572}{Louvain}{Profession of the four Jesuit vows.}{A;
exact civil day open}
\biotimelinerow{1576}{Rome}{Recalled to the Roman College chair of
Controversies.}{A/K strong}
\biotimelinerow{1576--1587/1588}{Rome}{Lectures on disputed Catholic and
Protestant claims.}{Institutional endpoint varies}
\biotimelinerow{1578}{Rome}{Publication of \emph{Institutiones linguae
Hebraicae}.}{Bibliographic year}
\biotimelinerow{1584}{Rome / print networks}{\emph{De translatione imperii}
appears as publication of controversial work develops.}{Bibliographic / K}
\biotimelinerow{1586--1593}{Ingolstadt}{Initial volumes of the
\emph{Disputationes} printed in stages.}{Early editions / K; revisions follow}
\biotimelinerow{1587}{Rome}{Begins service as consultor of the Congregation of
the Index.}{K; appointment act open}
\biotimelinerow{1588--1592}{Rome}{Spiritual father at the Roman College.}{K;
individual directees bounded}
\biotimelinerow{1589--1590}{France / Paris}{Papal legation and experience of
the siege of Paris.}{A/K; political reading interested}
\biotimelinerow{1590}{Rome}{Sixtus V moves to restrict the
\emph{Disputationes}.}{A/K; intended Index unpromulgated}
\biotimelinerow{1592}{Rome}{Work on correction of the Sixtine Vulgate and a
protective preface for the Sixto-Clementine edition.}{A/K; full collation open}
\biotimelinerow{1592--1595}{Rome}{Rector of the Roman College.}{A/K; one
chronology ends 1594}
\biotimelinerow{1595--1597}{Naples}{Jesuit provincial government.}{A/K strong}
\biotimelinerow{Jan.~1597}{Rome}{Summoned back to Roman service.}{A month/year}
\biotimelinerow{1597}{Rome}{Consultor of the Roman Holy Office.}{K; exact act
open}
\biotimelinerow{1597--1598}{Rome / print networks}{Brief and fuller catechisms
published.}{Bibliographic / M reception}
\biotimelinerow{14 Jan.~1599}{Rome}{Holy Office approves eight propositions
assembled from Bruno's works and dossier by Bellarmine and Tragagliolo.}{K from
surviving-summary reconstruction; list lost}
\biotimelinerow{3 Mar.~1599}{Rome}{Clement VIII creates Bellarmine cardinal.}{A/K
exact}
\biotimelinerow{17 Feb.~1600}{Rome}{Bruno is burned after the collective Roman
process.}{K exact outcome; responsibility not solitary}
\biotimelinerow{1600--1602}{Rome}{Participation and counsel in the
\emph{de auxiliis} deliberations.}{K broad sequence; memorials uncollated}
\biotimelinerow{Mar.--Apr.~1602}{Rome / Capua}{Election or appointment as
archbishop and episcopal consecration.}{K/M; sources give 18 or 25 March and 20
or 21 April for differently described stages; acts open}
\biotimelinerow{1 / 4 May 1602}{Capua}{Remembered arrival and formal entry to
begin residential government.}{A/K; private/formal distinction}
\biotimelinerow{1602--1605}{Capua}{Three diocesan synods, one provincial
council, visitation, preaching, and reported care for the poor.}{A/K/M;
outcomes incompletely audited}
\biotimelinerow{Mar. / Apr.--May 1605}{Rome}{Conclaves elect Leo XI and Paul V;
Bellarmine participates.}{A/K; near-election claim unresolved}
\biotimelinerow{1605}{Rome / Capua}{Returns to curial service and his Capuan
tenure ends.}{K sequence; resignation act and motive open}
\biotimelinerow{1606}{Rome / Venice}{Reported writings in the Venetian
Interdict controversy.}{K; exact writings and reciprocal dossier uncollated}
\biotimelinerow{1606--1609}{Rome / England by texts}{English Oath of Allegiance
controversy with Blackwell, James I, and others.}{A/K; no English journey}
\biotimelinerow{1608}{Ingolstadt}{Located printing of \emph{Recognitio librorum
omnium}.}{Bibliographic; claimed 1607 Rome printing unlocated}
\biotimelinerow{1610}{Rome}{\emph{De potestate Summi Pontificis in rebus
temporalibus} against Barclay.}{Bibliographic / A}
\biotimelinerow{1611}{Rome}{Publication of the Psalms commentary.}{Bibliographic
year}
\biotimelinerow{June 1613}{Rome}{Completion of the Latin autobiography.}{A
exact month; narrative then ends}
\biotimelinerow{12 Apr.~1615}{Rome}{Letter to Foscarini on Copernicanism,
demonstration, and scriptural interpretation.}{A/N exact document}
\biotimelinerow{1615--1620}{Rome}{Annual sequence of six late spiritual
books.}{Early editions / K}
\biotimelinerow{24--26 Feb.~1616}{Rome}{Assessment, papal direction,
Bellarmine's warning to Galileo, and the stronger precept record.}{N/K; record
relationship disputed}
\biotimelinerow{5 Mar.~1616}{Rome}{Index decree concerning Copernicus and
related books.}{N act; individual authorship not assumed}
\biotimelinerow{26 May 1616}{Rome}{Bellarmine certifies that Galileo had not
abjured or received penance while confirming the prohibition notified to
him.}{A/N exact document}
\biotimelinerow{1620}{Rome}{\emph{De arte bene moriendi}, final book in the
late annual sequence.}{Bibliographic year}
\biotimelinerow{10 Sept.~1621}{Rome}{Separately witnessed declaration of
adherence in the APUG manuscript addendum.}{N/E; not autobiography}
\biotimelinerow{17 Sept.~1621}{Sant'Andrea al Quirinale}{Death.}{N/E/K exact;
later official concurrence}
\biotimelinerow{1622}{Rome}{\emph{Adumbrata imago} and early cause-related
memory; exhumation and monument history begins.}{E/K; details need loci}
\biotimelinerow{1623--1624}{Rome}{Bernini bust and tomb monument produced
posthumously; Fuligatti's \emph{Vita} appears in 1624.}{E/K / object chronology}
\biotimelinerow{13 May 1923}{Rome}{Beatification by Pius XI.}{M exact}
\biotimelinerow{21 June 1923}{Rome}{Transfer of the remains then venerated from
the Gesù to Sant'Ignazio.}{M; not independent authentication}
\biotimelinerow{29 June 1930}{Vatican City}{Canonization by Pius XI; the
dispositive formula directs annual universal memory on 17 September.}{M exact}
\biotimelinerow{17 Sept.~1931}{Rome}{Declaration as Doctor of the Universal
Church; extension of the then 13 May Mass and Office.}{M exact historical act}
\biotimelinerow{29 Sept.~1967}{Capua}{Paul VI confirms and declares Bellarmine
equal principal patron of the Archdiocese with Saint Stephen Protomartyr.}{M
exact local act}
\biotimelinerow{20 July 2026 cutoff}{General Roman Calendar}{Optional memorial
on 17 September.}{M mutable; particular calendars may differ}
\end{biotimeline}

\subsection{Method, source limits, and negative results}

The narrative coordinates Bellarmine's autobiography and works, contemporary
institutional and reciprocal records, the first lives, official Catholic acts,
and named modern reconstruction.  An autobiographical motive is not converted
into an external fact.  An inquisitorial summary is not treated as a complete
tribunal dossier.  Cause-promoting testimony is read for proximity and genre.
Canonization and Doctor acts establish their own ecclesial judgments rather
than supplying neutral evidence for every sixteenth-century event.

The APUG autobiography transcription is marked ``To be proofread''; no
quotation depends on it as an unchecked authority.  Public-domain early
editions and the nineteenth-century Fèvre corpus are research witnesses, not
uniform modern critical editions.  The publication uses newly composed English
paraphrase and gives no detached quotation whose original locus, edition,
language, and translator remain uncontrolled.

No complete Bruno dossier or copy of the eight propositions was located.  The
procedural status of the unpromulgated Sixtine Index, full Vulgate revision
sequence, complete \emph{de auxiliis} file, Capuan administration, exact 1605
conclave ballots, reciprocal Venetian and English dossiers, and documentary
relationship of the 1616 Galileo warning and stronger precept remain open.
Those gaps do not make documented participation disappear and do not authorize
invented motives or votes.

No controlled source established an official lifetime title ``hammer of
heretics''; a journey to England for oath negotiation; membership in the
dicastery founded as Propaganda Fide after Bellarmine's death; exact early loci
for the popular sayings about walls, peace, or God's glory; a secure life
portrait; a complete physical custody chain or scientific authentication of
the present remains; or a competent universal patronage act for the frequently
claimed classes.  Paul~VI's bounded 1967 patronage act for the Archdiocese of
Capua is retained.  These are negative findings within the inspected corpus,
not proof that no other local or lost record exists.

\subsection{Rights and review state}

Project-created prose and organization are intended for CC BY 4.0.
Bellarmine's early-modern Latin and Italian and nineteenth-century editions are
public-domain textual candidates, but scans, manuscript images, modern
critical texts and apparatus, translations, scholarship, institutional web
presentation, and official-site prose retain their own rights and terms.  No
third-party image, facsimile, substantial modern translation, or protected
apparatus is reproduced.  Online access is not treated as a republication
license.

Internal source, chronology, tradition, and rights boundaries are recorded in
the publication and its research files.  Outstanding are independent
early-modern, Jesuit, theological, legal, institutional, opponent-side,
scientific-history, art-historical, relic, and ecclesiastical review; direct
collation of the complete major dossiers and editions; work-specific rights
review for distribution; and final production review.  The publication has no
imprimatur, nihil obstat, or ecclesiastical approval and remains a study aid.
